% GENERATED FILE. Edit canonical lessons, then run npm run generate:books.
% canonical-source-hash: fnv1a64:f5386612d82a7dff
% canonical-lessons: TA-C540-kacceri, TA-C540-natakam, TA-C540-metai, TA-C540-ottikai, TA-C540-oppanai

\chapter{On Stage}
\label{ch:ta-on-stage}
\hlchaptermodality{540}

\begin{chapteropening}
\textbf{By the end of this chapter:} \emph{I can talk about concerts, plays and rehearsals.}

Complete the last lesson of chapter 540: I can talk about concerts, plays and rehearsals.
\end{chapteropening}

\section[\texorpdfstring{\ta{கச்சேரி} (kaccēri)}{கச்சேரி (kaccēri)}]{\ta{கச்சேரி} (kaccēri) — a concert; a recital}

\label{lesson:TA-C540-kacceri}



\begin{quote}
Before the new one: say the Tamil for to pluck strings, then the Tamil for to compose music.
\end{quote}

\subsection*{You'll want to know: \ta{கச்சேரி}}
\textbf{\ta{கச்சேரி}} — \emph{kaccēri} — \textquotedblleft{}a concert; a recital\textquotedblright{}.

\begin{quote}
\ta{சென்ற} \ta{வாரம்} \ta{நாங்கள்} \ta{ஒரு} \ta{இசைக்} \ta{கச்சேரிக்குப்} \ta{போனோம்}.
\end{quote}

\begin{quote}
\emph{ceṉṟa vāram nāṅkaḷ oru icaik kaccērikkup pōṉōm.}
\end{quote}

\begin{quote}
Last week we went to a music concert.
\end{quote}

In South India it names a classical music recital.

\begin{cousinweb}
The origin of \textbf{\ta{கச்சேரி}} (\emph{kaccēri}) is not settled: the same word means an office or a court in other Indian languages, but the path into Tamil is unclear.
\end{cousinweb}

\begin{grammarlens}[title={What you've built}]
One more word for describing what you lived and felt.
\end{grammarlens}

\subsection*{Your turn}
\begin{itemize}
\raggedright
  \item \emph{Say it:} \emph{kaccēri}
  \item \emph{Say it:} \emph{kaccēri}, once more
  \item \emph{Say it:} \emph{ceṉṟa vāram nāṅkaḷ oru icaik kaccērikkup pōṉōm.} — Last week we went to a music concert.
  \item \emph{Recall:} say the Tamil for to pluck strings, then the Tamil for to compose music, then say \emph{kaccēri} again
\end{itemize}

\begin{tcolorbox}[breakable,colback=teal!4,colframe=teal!35!black,title={Before you move on}]
What does \ta{கச்சேரி} mean? (\textquotedblleft{}A concert; a recital\textquotedblright{}.) Say it once more.
\end{tcolorbox}
\section[\texorpdfstring{\ta{நாடகம்} (nāṭakam)}{நாடகம் (nāṭakam)}]{\ta{நாடகம்} (nāṭakam) — a stage play; a drama}

\label{lesson:TA-C540-natakam}



\begin{quote}
Before the new one: say the Tamil for to compose music, then the Tamil for a concert.
\end{quote}

\subsection*{You'll want to know: \ta{நாடகம்}}
\textbf{\ta{நாடகம்}} — \emph{nāṭakam} — \textquotedblleft{}a stage play; a drama\textquotedblright{}.

\begin{quote}
\ta{பள்ளி} \ta{ஆண்டு} \ta{விழாவில்} \ta{மாணவர்கள்} \ta{ஒரு} \ta{நாடகம்} \ta{நடித்தார்கள்}.
\end{quote}

\begin{quote}
\emph{paḷḷi āṇṭu viḻāvil māṇavarkaḷ oru nāṭakam naṭittārkaḷ.}
\end{quote}

\begin{quote}
At the school's annual day the students acted in a play.
\end{quote}

One acts in a play with \ta{நடி} (\emph{naṭi}) or stages one with \ta{போடு}: \ta{நாடகம்} \ta{போடு}.

\begin{cousinweb}
\textbf{\ta{நாடகம்}} (\emph{nāṭakam}) is borrowed from Sanskrit \emph{nāṭaka}, \textquotedblleft{}a dancer, an actor\textquotedblright{}, later \textquotedblleft{}a play\textquotedblright{} (CDIAL 7045).
\end{cousinweb}

\begin{grammarlens}[title={What you've built}]
One more word for describing what you lived and felt.
\end{grammarlens}

\subsection*{Your turn}
\begin{itemize}
\raggedright
  \item \emph{Say it:} \emph{nāṭakam}
  \item \emph{Say it:} \emph{nāṭakam}, once more
  \item \emph{Say it:} \emph{paḷḷi āṇṭu viḻāvil māṇavarkaḷ oru nāṭakam naṭittārkaḷ.} — At the school's annual day the students acted in a play.
  \item \emph{Recall:} say the Tamil for to compose music, then the Tamil for a concert, then say \emph{nāṭakam} again
\end{itemize}

\begin{tcolorbox}[breakable,colback=teal!4,colframe=teal!35!black,title={Before you move on}]
What does \ta{நாடகம்} mean? (\textquotedblleft{}A stage play; a drama\textquotedblright{}.) Say it once more.
\end{tcolorbox}
\section[\texorpdfstring{\ta{மேடை} (mēṭai)}{மேடை (mēṭai)}]{\ta{மேடை} (mēṭai) — a stage; a platform}

\label{lesson:TA-C540-metai}



\begin{quote}
Before the new one: say the Tamil for a concert, then the Tamil for a stage play.
\end{quote}

\subsection*{You'll want to know: \ta{மேடை}}
\textbf{\ta{மேடை}} — \emph{mēṭai} — \textquotedblleft{}a stage; a platform\textquotedblright{}.

\begin{quote}
\ta{பாடகி} \ta{மேடையில்} \ta{ஏறினாள்}.
\end{quote}

\begin{quote}
\emph{pāṭaki mēṭaiyil ēṟiṉāḷ.}
\end{quote}

\begin{quote}
The singer went up on the stage.
\end{quote}

Also a speaker's platform at a meeting. Where one stands takes -\ta{இல்}: \ta{மேடையில்}.

\begin{cousinweb}
The origin of \textbf{\ta{மேடை}} (\emph{mēṭai}) is not settled: no dictionary we can check gives its source.
\end{cousinweb}

\begin{grammarlens}[title={What you've built}]
One more word for describing what you lived and felt.
\end{grammarlens}

\subsection*{Your turn}
\begin{itemize}
\raggedright
  \item \emph{Say it:} \emph{mēṭai}
  \item \emph{Say it:} \emph{mēṭai}, once more
  \item \emph{Say it:} \emph{pāṭaki mēṭaiyil ēṟiṉāḷ.} — The singer went up on the stage.
  \item \emph{Recall:} say the Tamil for a concert, then the Tamil for a stage play, then say \emph{mēṭai} again
\end{itemize}

\begin{tcolorbox}[breakable,colback=teal!4,colframe=teal!35!black,title={Before you move on}]
What does \ta{மேடை} mean? (\textquotedblleft{}A stage; a platform\textquotedblright{}.) Say it once more.
\end{tcolorbox}
\section[\texorpdfstring{\ta{ஒத்திகை} (ottikai)}{ஒத்திகை (ottikai)}]{\ta{ஒத்திகை} (ottikai) — a rehearsal}

\label{lesson:TA-C540-ottikai}



\begin{quote}
Before the new one: say the Tamil for a stage play, then the Tamil for a stage.
\end{quote}

\subsection*{You'll want to know: \ta{ஒத்திகை}}
\textbf{\ta{ஒத்திகை}} — \emph{ottikai} — \textquotedblleft{}a rehearsal\textquotedblright{}.

\begin{quote}
\ta{நாளை} \ta{மாலை} \ta{கடைசி} \ta{ஒத்திகை} \ta{நடக்கும்}.
\end{quote}

\begin{quote}
\emph{nāḷai mālai kaṭaici ottikai naṭakkum.}
\end{quote}

\begin{quote}
The last rehearsal will take place tomorrow evening.
\end{quote}

With \ta{பார்} (\emph{pār}) it becomes a verb: \ta{ஒத்திகை} \ta{பார்}, to rehearse.

\begin{cousinweb}
The origin of \textbf{\ta{ஒத்திகை}} (\emph{ottikai}) is not settled: no dictionary we can check gives its source.
\end{cousinweb}

\begin{grammarlens}[title={What you've built}]
One more word for describing what you lived and felt.
\end{grammarlens}

\subsection*{Your turn}
\begin{itemize}
\raggedright
  \item \emph{Say it:} \emph{ottikai}
  \item \emph{Say it:} \emph{ottikai}, once more
  \item \emph{Say it:} \emph{nāḷai mālai kaṭaici ottikai naṭakkum.} — The last rehearsal will take place tomorrow evening.
  \item \emph{Recall:} say the Tamil for a stage play, then the Tamil for a stage, then say \emph{ottikai} again
\end{itemize}

\begin{tcolorbox}[breakable,colback=teal!4,colframe=teal!35!black,title={Before you move on}]
What does \ta{ஒத்திகை} mean? (\textquotedblleft{}A rehearsal\textquotedblright{}.) Say it once more.
\end{tcolorbox}
\section[\texorpdfstring{\ta{ஒப்பனை} (oppaṉai)}{ஒப்பனை (oppaṉai)}]{\ta{ஒப்பனை} (oppaṉai) — make-up; adornment}

\label{lesson:TA-C540-oppanai}



\begin{quote}
Before the new one: say the Tamil for a stage, then the Tamil for a rehearsal.
\end{quote}

\subsection*{You'll want to know: \ta{ஒப்பனை}}
\textbf{\ta{ஒப்பனை}} — \emph{oppaṉai} — \textquotedblleft{}make-up; adornment\textquotedblright{}.

\begin{quote}
\ta{நடிகர்கள்} \ta{நாடகத்துக்கு} \ta{முன்} \ta{ஒப்பனை} \ta{செய்தார்கள்}.
\end{quote}

\begin{quote}
\emph{naṭikarkaḷ nāṭakattukku muṉ oppaṉai ceytārkaḷ.}
\end{quote}

\begin{quote}
The actors put on make-up before the play.
\end{quote}

With \ta{செய்} (\emph{cey}): \ta{ஒப்பனை} \ta{செய்}, to put on make-up.

\begin{cousinweb}
\textbf{\ta{ஒப்பனை}} (\emph{oppaṉai}) is native Dravidian: DEDR 924, under \ta{ஒ} (\emph{o}), \textquotedblleft{}to suit, to resemble\textquotedblright{}, gives \ta{ஒப்பனை} (\emph{oppaṉai}), \textquotedblleft{}comparison, likeness, adornment\textquotedblright{}.
\end{cousinweb}

\begin{grammarlens}[title={What you've built}]
One more word for describing what you lived and felt.
\end{grammarlens}

\subsection*{Your turn}
\begin{itemize}
\raggedright
  \item \emph{Say it:} \emph{oppaṉai}
  \item \emph{Say it:} \emph{oppaṉai}, once more
  \item \emph{Say it:} \emph{naṭikarkaḷ nāṭakattukku muṉ oppaṉai ceytārkaḷ.} — The actors put on make-up before the play.
  \item \emph{Recall:} say the Tamil for a stage, then the Tamil for a rehearsal, then say \emph{oppaṉai} again
\end{itemize}

\begin{tcolorbox}[breakable,colback=teal!4,colframe=teal!35!black,title={Before you move on}]
What does \ta{ஒப்பனை} mean? (\textquotedblleft{}Make-up; adornment\textquotedblright{}.) Say it once more.
\end{tcolorbox}
